\chapter{Methodology}
\label{ch:2}

\lipsum[1]

\section{A table}
\label{sec:table}
\lipsum[1]


\begin{table}[htbp]
	\centering
	\captionof{table}{Constants of a system}
	\label{table:systemConstants}
	\begin{tabularx}{.75\textwidth}{XXX}
        \toprule
        Parameter & Value & Unit \\ \midrule \midrule
        $m$ & $0.19$ & $\text{kg}$ \\ \midrule
		$k$ & $0.2$ & $\frac{\text{kg}}{\text{s}}$ \\ \midrule
		$c$ & $4.8398$ & $\frac{\text{N}}{\text{m}}$ \\ \bottomrule
    \end{tabularx}
\end{table}

\section{Some equations}
\label{sec:someEquations}
\lipsum[1]


\begin{align*}
    \label{eq:system}
	\dt{\x} &= \A \x + \b u, &
	y &= \c^T \x + d u.
\end{align*}

\begin{equation}
	\label{eq:equation2}
	\begin{array}{c}
		\left[\begin{array}{c}
			 \x_{k+1} \\
			 u_{k}
			\end{array}
		\right]
		 = 
		\underbrace{
			\left[ 
				\begin{array}{cc}
					\A_d & \mt{\Gamma}_1(\tau_a, \tau_b) \\
					\mt{0} &  0 \\
				\end{array}
			\right]
		}_{=: \A_{\xi}(\tau_a, \tau_b)}
		\underbrace{
			\left[ 
				\begin{array}{c}
					\x_{k} \\
					u_{k-1}
				\end{array}
			\right]
		}_{=: \vc{\xi}_{k}} 
		+
		\underbrace{
			\left[
				\begin{array}{c}
					\mt{\Gamma}_0(\tau_a, \tau_b) \\
					1
				\end{array}
			\right]
		}_{=: \b_{\xi}(\tau_a, \tau_b)} u_k ,
		\\ \\
		y_k = \left[ \c^T \quad 0 \right] \vc{\xi}_k ,
	\end{array}
\end{equation}


\begin{equation}
	\label{eq:discSystemLargeDelayCompact}
	\begin{array}{c}
		\left[\begin{array}{c}
			\x_{k+1} \\
			u_{k} \\
			u_{k-1} \\
			\vdots \\
			u_{k-\zeta+1} \\
			\end{array}
		\right]
		= 
		\underbrace{\left[ \begin{array}{ccccc}
			\A_d & \Gamma_{1}(\mathrm{Y}) & \hdots & \Gamma_{\zeta - 1}(\mathrm{Y}) & \Gamma_{\zeta}(\mathrm{Y}) \\
			\mt{0} & 0 & \hdots & 0 & 0 \\
			\mt{0} & 1 & \hdots & 0 & 0 \\
			\vdots &  \vdots  & \ddots & \vdots & \vdots \\
			\mt{0} & 0 & \hdots & 1 & 0 \\
			
		\end{array}
		\right]}_{=: \A_{\xi}(\mathrm{Y})}
		\underbrace{
			\left[\begin{array}{c}
				\x_{k} \\
				u_{k-1} \\
				u_{k-2} \\
				\vdots \\
				u_{k-\zeta} \\
				\end{array}
			\right]
		}_{=: \vc{\xi}_{k}} \\ \\ \qquad \qquad 
		+ 
		\underbrace{
			\left[
				\begin{array}{c}
					\Gamma_{0}(\mathrm{Y}) \\
					1 \\
					0 \\
					\vdots \\
					0
				\end{array}
			\right]
		}_{=: \vc{b}_{\xi}(\mathrm{Y})} u_k ,
		\\ \\
		y_k = \left[ \c^T \quad 0 \quad 0 \quad \hdots \quad 0\right] \vc{\xi}_k .
	\end{array}
\end{equation}


\newpage
\section{Theorem}
\label{sec:theorem}

\subsection{Theorem 1}

\begin{theo}[Pythagoras' theorem]{thm:pythagoras}
	In a right triangle, the square of the hypotenuse is equal to the sum of the squares of the catheti.
	\[a^2+b^2=c^2\]
\end{theo}
In mathematics, the Pythagorean theorem, also known as Pythagoras' theorem, is a relation in Euclidean geometry among the three sides of a right triangle.


\subsection{Theorem 2 with proof}
\begin{theo}{thm:theorem1}
	There exist two irrational numbers $x$, $y$ such that $x^y$ is rational.
\end{theo}
 
\begin{prf}{prf:proof1}
	If $x=y=\sqrt{2}$ is an example, then we are done; otherwise $\sqrt{2}^{\sqrt{2}}$ is irrational, in which case taking $x=\sqrt{2}^{\sqrt{2}}$ and $y=\sqrt{2}$ gives us:
	\[\bigg(\sqrt{2}^{\sqrt{2}}\bigg)^{\sqrt{2}}=\sqrt{2}^{\sqrt{2}\sqrt{2}}=\sqrt{2}^{2}=2.\]
\end{prf}